\documentclass[10pt]{article}

\usepackage[utf8]{inputenc}

\usepackage[T1]{fontenc}

\usepackage{amsmath}

\usepackage{amsfonts}

\usepackage{amssymb}

\usepackage[version=4]{mhchem}

\usepackage{stmaryrd}



\begin{document}

число решений для любого комплексного числа $a \neq 0$. Эти решения находятся по формуле



\[

z=\operatorname{Ln} a=\ln |a|+i \operatorname{Arg} a .

\]



П. ф. $e^{z}$ является одной из основных әлементарных функций. Через нее выражаются, напр., тригонометрич. функции и гиперболич. функции. ГО. В. Сидоров.



ПОКАЗАТЕЛЬНОЕ РАСППЕДЕЛЕНИЕ - НеПрерывное распределение вероятностей случайной величины $X$, задаваемое плотностью



\[

p(x)= \begin{cases}\lambda e^{-\lambda x}, & x \geqslant 0, \\ 0, & x<0 .\end{cases}

\]



Плотность $p(x)$ зависит от положительного масштабного параметра $\lambda$. Формула для моментов: $E X^{n}=n ! / \lambda^{n}$, в частности - для математич. ожидания $\mathrm{E} X=1 / \lambda$ и дисперсии $\mathrm{D} X=1 / \lambda^{2}$; характеристич. фуннция: $(1-i t / \lambda)^{-1}$.



П. р. входит в семейство распределений, называемых гамма-распределениями и задаваемых плотностью



\[

p(x)=\left[\lambda^{\alpha} x^{\alpha-1} / \Gamma(\alpha)\right] e^{-\lambda x}, x \geqslant 0 ;

\]



$n$-кратная свертка распределения (1) равна гаммараспределению с тем же самым параметром $\lambda$ и с $\alpha=n$.



П. p.- единственное распределение, обладающее свойством отсутствия последействия: для любых $x>0$, $y>0$ выполняется равенство



\[

\mathbf{P}\{X>x+y \mid X>y\}=\mathbf{P}\{X>x\},

\]



где $\mathrm{P}\{X>x+y \mid X>y\}-$ условная вероятность события $X>x+y$ при условии $X>y$. Свойство (2) называется также м а р к о в с к и м с в ой с т в м.



В однородном пуассоновском процессе расстояние между двумя последовательными скачками траектории имеет П. р. Наоборот, процесс восстановления с показательным временем жизни (1) является пуассоновским процессом восстановления. П. р. часто возникает как предельное при суперпозиции или разрежевии процессов восстановления, в задачах пересечения высокого уровня в различных схемах блуждания, в критических ветвящихся процессах и т. п.



Упомянутыми выше свойствами объясняется широкое применение П. р. при расчетах различных систем в теории массового обслуживания и в теории надежности. Предполагая времена занятости приборов случайными, независимыми друг от друга и распределенными показательно, можно благодаря свойству (2) изучать системы массового обслуживания с помощью конечных или счетных цепей Маркова с непрерывным временем. Аналогичным образом используются цепи Маркова и в теории надежности, где времена исправной работы отдельных приборов часто можно предполагать независимыми и распределенными показательно.



Лит.: [1] Фе л л е р В., Введение в теорию вероятностей и ее приложения, пер. с англ., 2 изд., т. 2, М., 1967.



покоординатного СПтска Б. А. Севаствянов. ПОКООРДИНАТНОГО СПУСКА МЕТОД - одиН Из методов минимизации функций многих переменных, использующий лишь значения минимизируемой функции. П. с. м. применяется в тех случаях, когда минимизируемая функция недифференцируема или вычисление ее производных требует большого объема работы. Ниже описан П. с. м. для задачи минимизации функции $F(x)$ на множестве



\[

X=\left\{x=\left(x^{1}, \ldots, x^{n}\right): a_{i} \leqslant x^{i} \leqslant b_{i}, i=1, \ldots, n\right\},

\]



где $a_{i}, b_{i}$ - заданные числа, $a_{i}<b_{i}$; случаи, когда все или нек-рые $a_{i}=-\infty, b_{j}=+\infty$, здесь не исключаются. Пусть $e_{i}=(0, \ldots, 0,1,0, \ldots, 0)-$ координатный вектор, у к-рого $i$-я координата равна 1 , остальные координаты равны нулю. Задают начальное приближение $x_{0} \in X, \alpha_{0}>0$. Пусть известно $k$-е при- ближение $x_{0} \in X, \alpha_{k}>0$ при каком-либо $k \geqslant 0$. Полагают $\boldsymbol{p}_{k}=\boldsymbol{e}_{i_{k}}$, где $i_{k}=k-n\left[\frac{k}{n}\right]+1$ (здесь $[a]-$ целая часть числа a). Таким образом,



$\boldsymbol{p}_{0}=\boldsymbol{e}_{1}, \ldots, \boldsymbol{p}_{n-1}=\boldsymbol{e}_{n}, \boldsymbol{p}_{n}=\boldsymbol{e}_{1}, \ldots, \boldsymbol{p}_{2 n-1}=\boldsymbol{e}_{n}, \boldsymbol{p}_{2 n}=\boldsymbol{e}_{1}, \ldots$, т. е. осуществляется циклич. перебор координатных векторов $e_{1}, \ldots, e_{n}$. Сначала проверяют выполнение условия



\[

x_{k}+\alpha_{k} \boldsymbol{p}_{k} \in X, F\left(\boldsymbol{x}_{k}+\alpha_{k} \boldsymbol{p}_{k}\right)<F\left(\boldsymbol{x}_{k}\right) .

\]



Если (1) выполняется, то полагают $x_{k+1}=x_{k}+\alpha_{k} p_{k}$, $\alpha_{k+1}=\alpha_{k}$. Если (1) не выполняется, то проверяют условие



\[

\boldsymbol{x}_{k}-\alpha_{k} \boldsymbol{p}_{k} \in X, F\left(\boldsymbol{x}_{k}-\alpha_{k} \boldsymbol{p}_{k}\right)<F\left(\boldsymbol{x}_{k}\right) .

\]



В случае выполнения условия (2) полагают $x_{k+1}=$ $=x_{k}-\alpha_{k} p_{k}, \alpha_{k+1}=\alpha_{k}$. Если оба условия (1), (2) не выполняются, то полагают $\boldsymbol{x}_{k+1}=\boldsymbol{x}_{k}$,



$\alpha_{k+1}=\left\{\begin{array}{l}\lambda \alpha_{k} \text { при } i_{k}=n, x_{k}=x_{k-n+1}, \\ \alpha_{k} \text { при } i_{k} \neq n \text { или } x_{k} \neq x_{k-n+1}, \text { или } 0 \leqslant k \leqslant n-1 ;\end{array}\right.$



где $\lambda, 0<\lambda<1,-$ параметр метода. Условия (3) означают, что если за один цикл из $n$ итераций при переборе всех координатных векторов $e_{1}, \ldots, e_{n}$ с шагом $\alpha_{k}$ выполнилось хотя бы одно из условий (1) или (2), то длина шага $\alpha_{k}$ не дробится и сохраняется на протяжении по крайней мере следующего цикла из $n$ итераций; если же на последних $n$ итерациях оба условия (1), (2) ни разу не выполнились, то шаг $\alpha_{k}$ дробится.



Если функция $F(x)$ выпукла и непрерывно дифференцируема на $X$, множество $\left\{x \in X: F(x) \leqslant F\left(x_{0}\right)\right\}$ ограничено, $\alpha_{0}$ - произвольное положительное число, то метод (1) - (3) сходится, т. е.



\[

\lim _{k \rightarrow \infty} F\left(x_{k}\right)=\inf _{x \in X} F(x),

\]



последовательность $\left\{x_{k}\right\}$ сходится к множеству точек минимума $F(x)$ на $X$. Если $F(x)$ недифференцируема на $X$, то П. с. м. может не сходиться (см. [1], [2]).



Лит.: [1] В а с и ль е в Ф. П., Численные методы решения экстремальных зарач, M., 1980;' [2] K a p M a о o B B. Г., Математическое программирование, 2 изд., М., 1980.



ПОКРЫВАЮЩИИ ЭЛЕМЕНТ в ч а с Т И ч н 0 у по р я д о ч е н но м множе с т в е - элемент, непосредственно следующий за другим элементом; точнее, выражение " $a$ покрывает $b$ в частично упорядоченном множестве $P_{\text {》 означает, что }}^{b<a}$ и не существует элемента $x \in P$, удовлетворяющего условию $b<$ $<x<a$.



Т. C. Фофанова.



ПОКРЫТИЕ м н о ж с т в а $X-$ любое семейство подмножеств этого множества, объединение к-рого есть $X$.



\begin{enumerate}

  \item Под П. то пологическиго пр ос т р а н с т в а, равномерного пространства и вообще какого-либо мнонества, наделенного тем или иным строением, понимают произвольное П. этого множества. Однако в теории топологич. пространств особенно естественно рассматривать о т к р ы т ы е пок р ы т и я, то есть П., все элементы к-рых являются открытыми множествами. Большое значение открытых П. вызвано тем, что их элементы несут в себе полную информацию о локальном строении пространства, а свойства П. в целом (в частности, число элементов в пем, кратность, комбинаторные свойства) отражают существенно глобальные характеристики пространства.

\end{enumerate}



Так, на языке открытых П. определяется размерность по Јебегу $\operatorname{dim}$ топологич. пространс1ва: размерность нормального пространства $X$ не превосходит





\end{document}